% Clinical Background section. Paste into your own thesis, or \input it.
%
% FIGURE WANTED: a stage diagram of atherosclerosis -- healthy wall, fatty streak,
% fibrous cap, rupture with thrombus. Korona fig. 3.4 uses one and it carries the
% section. Not added yet: an external figure needs a licensed credit generated into
% figures/external/CREDITS.md by scripts/fetch_external_figures.py, which does not
% exist on this branch. Do not add a \ref until the figure and its credit exist.
% Clinical Background section. Paste into your own thesis, or \input it.
%
% FIGURE WANTED: a stage diagram of atherosclerosis -- healthy wall, fatty streak,
% fibrous cap, rupture with thrombus. Korona fig. 3.4 uses one and it carries the
% section. Not added yet: an external figure needs a licensed credit generated into
% figures/external/CREDITS.md by scripts/fetch_external_figures.py, which does not
% exist on this branch. Do not add a \ref until the figure and its credit exist.

\section{Coronary artery disease}
\label{sec:cad}

Cardiovascular disease is the leading cause of death worldwide and accounted for approximately
20.5 million deaths in 2021 \cite{lindstrom2022global}.
\textbf{Coronary artery disease (CAD)} is its largest single component.
It is what happens when \emph{atherosclerosis} develops in the arteries that supply the heart
muscle.

Atherosclerosis is slow.
Over years the inner wall of an artery takes up fat and calcium, and the deposit that builds up
there is called a \emph{plaque} \cite{libby2019atherosclerosis}.
As the deposit grows, a cap of fibrous tissue forms over it and walls it off from the blood
\cite{libby2019atherosclerosis}.
The plaque bulges inward as it enlarges and narrows the \emph{lumen}, the open channel the blood
flows through, so less blood can get past \cite{libby2019atherosclerosis}.

The coronary arteries are the dedicated supply of the myocardium (Section~\ref{sec:anatomy}), so a plaque in one of them is dangerous in two ways, which differ in speed. 
The first is gradual. Most tissues can take extra oxygen out of the blood already reaching them when demand rises, but
the heart cannot: it is already taking 60--80\% of what arrives, even at rest
\cite{duncker2015regulation}. Its only option is to receive more blood, and the plaque is what prevents that. The muscle beyond the plaque is left short of oxygen whenever the heart works harder, a state
called \emph{ischemia}.
The second is abrupt and far more dangerous. The fibrous cap tears, which
exposes the material inside the plaque to the bloodstream; blood clots on contact with that
material, and the clot, or thrombus, occludes the vessel \cite{rampidis2022geometry}. The muscle beyond the blockage is cut off from oxygen and dies within minutes. That is a myocardial infarction, also known as a heart attack.

All of this happens quietly.
CAD can advance for years without causing a single symptom \cite{knuuti2019esc}, and by the time symptoms do appear the disease is usually well established.
For some people the first sign of it is the heart attack.
Waiting for symptoms therefore means waiting too long.
The patients worth finding are the ones who still feel completely well, and finding them means
having something measurable that marks them out in advance.

The shape of the arteries is a candidate.
Plaque does not form just anywhere: the \emph{endothelium}, the layer of cells lining the artery,
reacts to how blood flows over it and turns inflamed where that flow is slow or unsteady
\cite{rampidis2022geometry}.
Those conditions arise at bends and where an artery branches, which are properties of the shape of
the tree (Sections~\ref{sec:hemodynamics} and~\ref{sec:geometry}).
Unlike a symptom, that shape can be measured on a scan of someone who feels perfectly healthy, and
Section~\ref{sec:ccta} describes the scan used to do it.

